% Generated draft: requires mathematical and source review.
\documentclass[11pt]{article}
\usepackage[margin=25mm]{geometry}
\usepackage{fontspec}
\setmainfont{Latin Modern Roman}
\usepackage{amsmath,amssymb,mathtools,unicode-math}
\setmathfont{Latin Modern Math}
\usepackage{enumitem,needspace,xurl,hyperref}
\setlength{\parindent}{0pt}
\setlength{\parskip}{5pt}
\setlength{\emergencystretch}{3em}
\newcommand{\N}{\mathbb{N}}
\newcommand{\Z}{\mathbb{Z}}
\newcommand{\Q}{\mathbb{Q}}
\newcommand{\R}{\mathbb{R}}
\newcommand{\C}{\mathbb{C}}
\newcommand{\F}{\mathbb{F}}
\newcommand{\A}{\mathbb{A}}
\newcommand{\PP}{\mathbb{P}}
\newcommand{\E}{\mathbb{E}}
\newcommand{\eps}{\varepsilon}
\newcommand{\dd}{\,\mathrm{d}}
\begin{document}
\section*{UnsolvedMath ID: 800429}
\textbf{Source problem code:} catalogue entry
\subsection{Definitions and mathematical statement}
Let \(N\geq 1\) denote the projective dimension, as implicitly intended in the supplied statement, and let
\[
f:\mathbb{P}^N_{\mathbb{Q}}\longrightarrow\mathbb{P}^N_{\mathbb{Q}}
\]
be a morphism of degree \(d\geq 2\). Write \(f^0=\mathrm{id}\), and let \(f^n\) denote the \(n\)-fold iterate. For each \(n\geq 1\), choose a homogeneous presentation
\[
f^n=[F_0^{(n)}:\cdots:F_N^{(n)}],
\qquad F_i^{(n)}\in\mathbb{Q}[X_0,\ldots,X_N],
\]
whose coordinate forms have degree \(d^n\) and no common projective zero. Assume that
\[
F_0^{(n)}\neq cX_0^{d^n}
\quad\text{for every }c\in\mathbb{Q}
\quad\text{and every }1\leq n\leq N+1.
\]
This condition is independent of the common nonzero scalar used in the homogeneous presentation.

Set \(H=\{X_0=0\}\). A rational point \(Q\notin H\) is integral relative to \(H\) in the standard coordinates if
\[
Q=[1:b_1:\cdots:b_N],\qquad b_1,\ldots,b_N\in\mathbb{Z}.
\]
Equivalently, if \(Q=[a_0:\cdots:a_N]\) is represented by relatively prime integer coordinates, then \(Q\) is integral relative to \(H\) exactly when \(|a_0|=1\). Define
\[
\mathcal{O}_f(P)=\{f^n(P):n\geq 0\}.
\]
A point \(P\) is wandering if this orbit is infinite.

\textbf{Supplied conjectural assertion (Conjecture 52).} For every \(P\in\mathbb{P}^N(\mathbb{Q})\),
\[
\#\bigl\{Q\in\mathcal{O}_f(P):Q\text{ is integral relative to }H\bigr\}<\infty.
\]
The source further gives the following assertion as a more precise formulation. If \(P\) is wandering, choose, for every \(n\geq 1\), coordinates
\[
f^n(P)=[A_0(n):\cdots:A_N(n)],\qquad
A_i(n)\in\mathbb{Z},\qquad
\gcd(A_0(n),\ldots,A_N(n))=1.
\]
Then the asserted limit is
\[
\lim_{n\to\infty}
\frac{\log|A_0(n)|}
{\log\max_{0\leq i\leq N}|A_i(n)|}=1.
\]
The absolute values make this expression independent of the simultaneous sign choice for the primitive coordinates.

\textbf{Unresolved qualification.} The displayed quotient requires \(A_0(n)\neq 0\) and \(\max_i|A_i(n)|>1\). A real-valued limit therefore requires these conditions on an eventual tail. The supplied statement does not specify eventual avoidance of \(H\), a convention for \(\log 0\), or a restriction to selected indices. It also states no Zariski-density assumption on the orbit and no exclusion of orbits in proper invariant subvarieties. Whether an additional orbit hypothesis is intended requires source and mathematical review; none has been inserted into the assertions above.
\subsection{Short English statement}
The supplied Dynamical Lang--Siegel conjecture predicts that a rational orbit contains only finitely many points with integral affine coordinates relative to \(X_0=0\), under the stated restriction on the first \(N+1\) iterates. For a wandering point it further asserts that the logarithm of the first primitive coordinate is asymptotic to the logarithmic projective height. Possible missing orbit hypotheses and the meaning of this limit when the first coordinate vanishes remain unresolved.
\subsection{Sources}
\begin{itemize}
\item Local UnsolvedMath registry, supplied complete record: permanent ID \(800429\), exact source problem code \texttt{catalogue entry}; title: \emph{Open Problems in Arithmetic Dynamics --- Conjecture 52 (Dynamical Lang--Siegel Conjecture for \(\mathbb{P}^N\))}. The supplied registry labels the problem \texttt{open}; this is a source status label, not an independently verified literature assessment.
\item Benedetto, DeMarco, Ingram, Jones et al., \emph{Current Trends and Open Problems in Arithmetic Dynamics} (2018), Conjecture 52. Supplied reference: \url{https://arxiv.org/abs/1806.04980}. Attribution and proposed year are as supplied in the registry.
\item Supplied provenance records access on 2026-10-16 and extraction by \texttt{arxiv\_latex}. These are supplied provenance claims; the reference and extracted formulation have not been independently checked for this draft. The registry records creation and update on 2026-10-16, with upstream revision \texttt{d371ee917ac9647fbb06cd6fc529deec6f3dda82}.
\end{itemize}
\end{document}
